%% LyX 2.5.1 created this file.  For more info, see https://www.lyx.org/.
%% Do not edit unless you really know what you are doing.
\documentclass[english,hebrew]{article}
\usepackage[HE8,T1]{fontenc}
\usepackage[utf8]{inputenc}
\usepackage{babel}
\usepackage{amsmath}
\usepackage{amssymb}
\usepackage[unicode=true]
 {hyperref}
\begin{document}
\title{מבוא לטופולוגיה קבוצתית: קומפקטיות}
\maketitle
\begin{description}
\item [{קטגוריות:}] טופולוגיה
\item [{תגים:}] מרחבים טופולוגיים, מרחבי מכפלה
\item [{מזהה:}] \L{topological\_spaces\_compactness}
\end{description}
סדרת הפוסטים שלי על טופולוגיה קבוצתית הגיעה לשלב שבו קצת פחות ברור
לאן צריך ללכת. ראינו את מושגי הבסיס המרכזיים: טופולוגיה, קבוצות פתוחות
וסגורות, נקודות הצטברות, פונקציות רציפות, טופולוגיית הסדר, טופולוגיית
תת-המרחב, טופולוגיית המכפלה... מכאן אפשר ללכת לשלל כיוונים שונים ומשונים,
אז אני הולך אל האחד שאני אוהב במיוחד אישית: מושג \textbf{הקומפקטיות}.
כאן יש לנו המחשה שלדעתי עובדת מצוין לאופן שבו טופולוגיה קבוצתית לא
סתם מכלילה רעיונות בסיסיים בצורה חכמה - היא גם עושה את זה \textbf{ממש
מוזר}. בניסוח הבסיסי שלה, קומפקטיות נשמעת כמו תכונה תלושה לחלוטין
ולא כמו הדבר החזק והמועיל שהיא. ומה הניסוח הזה? אוקיי, זה קל:
\begin{itemize}
\item מרחב טופולוגי \L{$X$} הוא \textbf{קומפקטי} אם לכל כיסוי פתוח שלו
יש תת-כיסוי סופי. כלומר, פורמלית: לכל אוסף \L{$\left\{ U_{\alpha}\right\} _{\alpha\in\Lambda}$}
של קבוצות פתוחות ב-\L{$X$} כך ש-\L{$\bigcup_{\alpha\in\Lambda}U_{\alpha}=X$}
קיים תת-אוסף \L{$\Lambda^{\prime}\subseteq\Lambda$} \textbf{סופי}
כך ש-\L{$\bigcup_{\alpha\in\Lambda^{\prime}}U_{\alpha}=X$}
\end{itemize}
אוקיי... מה?

{*}חריקת תקליט{*}

{*}הקפאת תמונה{*}

\textquotedblleft כן, זה אני. אתם בוודאי תוהים איך הגעתי למצב הזה.
הכל התחיל לפני כמה שבועות...\textquotedblright{}

\section{כמה שבועות קודם לכן, הישר הממשי}

אחד הדברים הראשונים שאפשר לדבר עליהם בחשבון אינפיניטסימלי אחרי שיש
לנו את המושג של פונקציות רציפות הוא \textbf{משפטי ויירשטראס}. יש שני
משפטים שאפשר לאחד עם ניסוח מספיק זהיר:
\begin{itemize}
\item פונקציה רציפה על קטע סגור וחסום \L{$f:\left[a.b\right]\to\mathbb{R}$}
חסומה בקטע הזה.
\item פונקציה רציפה על קטע סגור וחסום \L{$f:\left[a.b\right]\to\mathbb{R}$}
מקבלת את המקסימום והמינימום שלה בקטע הזה.
\end{itemize}
אפשר לאחד אותם אל הטענה הבאה: בהינתן פונקציה רציפה על קטע סגור וחסום
\L{$f:\left[a.b\right]\to\mathbb{R}$}, קיימות שתי נקודות \L{$c,C\in\left[a,b\right]$}
כך ש-\L{$f\left(c\right)\le f\left(x\right)\le f\left(C\right)$}
לכל \L{$x\in\left[a,b\right]$}. הסיבה שיש שני משפטים היא שנוח קונספטואלית
להוכיח קודם את הראשון בתור תוצאה נפרדת ואז לקבל את השני בתור המשך
טבעי שלו. אני לא אכנס עכשיו לשימושיות של המשפטים הללו אבל הם מאוד
שימושיים; בפרט, הם מהווים את הצעד הראשון בדרך להגדרה של \textbf{נגזרות}.

איך מוכיחים את זה? הנה הוכחה שיש לי כרגע בראש. עבור המשפט הראשון,
בואו נניח בשלילה ש-\L{$f$} לא חסומה. אז לכל \L{$n$} טבעי, קיים \L{$x_{n}\in\left[a,b\right]$}
כך ש-\L{$f\left(x_{n}\right)>n$}. קיבלנו סדרה \L{$\left\{ x_{n}\right\} ^{\infty}_{n=1}$}
של נקודות שכולן נמצאות בקטע \textbf{החסום} \L{$\left[a,b\right]$},
אז לפי \textbf{משפט בולצאנו-ויירשטראס} יש להן תת-סדרה מתכנסת \L{$x_{n_{i}}$}.
נסמן את הגבול שלה ב-\L{$x=\lim_{i\to\infty}x_{n_{i}}$}. בגלל שהקטע
\L{$\left[a,b\right]$} \textbf{סגור} אז \L{$x\in\left[a,b\right]$}
גם כן. עכשיו, נשתמש בתכונה של פונקציות ממשיות רציפות: \L{$\lim_{i\to\infty}f\left(x_{n_{i}}\right)=f\left(\lim_{i\to\infty}x_{n_{i}}\right)=f\left(x\right)$},
אלא שלרוע המזל, \L{$\lim_{i\to\infty}f\left(x_{n_{i}}\right)=\infty$}
פשוט מההגדרה, כי לכל \L{$N$} אפשר למצוא מקום בסדרה שהחל ממנו \L{$f\left(x_{n_{i}}\right)>N$}
כפי שקיבלנו מההנחה ש-\L{$f$} לא חסומה. מכיוון שלא ייתכן ש-\L{$f\left(x\right)=\infty$}
)אין לזה משמעות פורמלית בכלל(, הגענו לסתירה ולכן \L{$f$} כן חסומה.

בשביל המשפט השני אפשר עכשיו להניח שכבר ידוע ש-\L{$f$} חסומה. נסתכל
על קבוצת כל הערכים שלה, \L{$f\left(\left[a,b\right]\right)$}. הקבוצה
הזו חסומה מלעיל, ולכן יש לה \textbf{חסם עליון} \L{$M$} כתוצאה מ\textbf{אקסיומת
השלמות} של המספרים הממשיים. המשמעות של חסם עליון היא שלכל \L{$\varepsilon>0$},
קיים \L{$x\in\left[a,b\right]$} כך ש-\L{$f\left(x\right)>M-\varepsilon$}.
אז נמצא \L{$x_{n}$} כזה לכל \L{$\varepsilon=\frac{1}{n}$} וקיבלנו
סדרה \L{$\left\{ x_{n}\right\} ^{\infty}_{n=1}$} כך ש-\L{$f\left(x_{n}\right)\to M$}.
כמו קודם, ניקח תת-סדרה מתכנסת של הסדרה הזו עם גבול \L{$x$} ואז מרציפות
\L{$f$} יתקיים ש-\L{$f\left(x\right)=M$}. באותה דרך מראים שגם המינימום
של \L{$f$} מתקבל.

זה מסיים את ההוכחות, שהן לדעתי פחות או יותר הדבר הסטנדרטי שרואים כשלומדים
חשבון אינפיניטסימלי. ההסתייגות היחידה שלי היא שיש כאלו שמלמדים אינפי
ומתחילים ישר מפונקציות בלי לדבר על סדרות ועל משפט בולצאנו ויירשטראס,
מה שדי מסרבל את ההוכחות הללו ועוד כמה )יש לי עדיין טראומה מההוכחה
המסורבלת הרבה יותר מדי של משפט קנטור על רציפות במידה שווה שהיה בספר
קורס האינפי שלמדתי באוניברסיטה הפתוחה(.

עכשיו, בטופולוגיה מה שאנחנו עושים הוא לקחת דברים מגניבים מאינפי ולהכליל
אותם. איך נעשה את זה עם משפטי ויישרטאס? קודם כל בואו נסכם באילו רכיבים
השתמשנו פה:
\begin{itemize}
\item \textbf{אקסיומת השלמות}
\item משפט\textbf{ בולצאנו ויירשטראס}
\item פונקציה \textbf{רציפה}
\item קטע \textbf{חסום}
\item קטע \textbf{סגור}
\end{itemize}
על אקסיומת השלמות ובולצאנו ויירשטראס אני מדבר \href{https://gadial.net/2024/10/12/real_numbers_two_completness/}{באחד מהפוסטים}
בסדרת הפוסטים שלי על ההגדרה של \textbf{המספרים הממשיים} )וגם שם אני
מוכיח את משפטי ויירשטראס(. בגדול, מבלי להיכנס לפרטים, אקסיומת השלמות
היא התכונה המהותית ביותר של המספרים הממשיים, ובולצאנו ויירשטראס הוא
תוצאה מיידית ומתבקשת שלה. אז שני המושגים הללו הם אפיון מיוחד של \L{$\mathbb{R}$}
ולכן אם רוצים להכליל את המשפטים הללו צריך איכשהו לשמר את ההיבטים המהותיים
שלהם שבהם משתמשים בהוכחות הללו. אם היינו מעל \L{$\mathbb{Q}$}, למשל,
המשפטים פשוט לא היו עובדים. יש דוגמא מעליבה ממש, \L{$f\left(x\right)=\frac{1}{x^{2}-2}$}
מעל הקטע \L{$\left[1,2\right]$} - הפונקציה הזו רציפה בכל מספר \textbf{רציונלי},
כי הנקודה היחידה שבעייתית היא זו שבה המכנה מתאפס - וזה \L{$x=\sqrt{2}$}
האי-רציונלי. מצד שני, הפונקציה הזו מאוד לא חסומה כי ככל שלוקחים מספרים
רציונליים שהולכים ומתקרבים אל \L{$\sqrt{2}$}, הערך של הפונקציה \textquotedblleft מתפוצץ\textquotedblright .
אז כל נסיון הכללה של המשפט ידרוש להשתמש בתכונה של \L{$\mathbb{R}$}
שאין ב-\L{$\mathbb{Q}$} ולא ברור מה התכונה הזו צריכה להיות.

וחוץ מזה, יש גם את העניין הזה עם \textquotedblleft קטע סגור וחסום\textquotedblright .
את \textbf{סגור} יש לנו, אבל מה עם חסום? בואו לא נתייאש, אפשר להגדיר
\textbf{חסום} במרחבים מטריים כלליים. אז בואו נתחיל מהם.

\section{מרחב מטרי, שעת לילה מאוחרת}

אם \L{$\left(X,d\right)$} הוא מרחב מטרי, הוא עדיין מזכיר את \L{$\mathbb{R}$}
בשלל דרכים אבל דבר אחד מהותי הולך מייד לאיבוד - אין על המרחב הזה \textbf{יחס
סדר}. אם אין יחס סדר אין בכלל משמעות לדברים כמו \textquotedblleft מקסימום\textquotedblright .
אז מראש אני לא הולך להכליל את המשפט \textbf{לכל} מרחב טופולוגי, אבל
מצד שני - אני צריך מקסימום ומינימום רק \textbf{בטווח} של הפונקציה,
התחום יכול להיות כל דבר. אז בשלב ראשון נסתכל על פונקציות \L{$f:\left(X,d\right)\to\mathbb{R}$}
ונשאל את עצמנו מה המרחב המטרי \L{$X$} צריך לקיים כדי שעדיין יהיה
לנו את משפטי ויירשטראס.

ממה שראינו קודם, נראה די ברור ש-\L{$X$} אמור להיות סגור, לא? אבל
הנה העניין, \L{$X$} הוא תמיד סגור - הוא כל המרחב! במשפט ויירשטראס
הקלאסי, נקודת המבט שלנו הייתה \textquotedblleft העולם שלנו הוא \L{$\mathbb{R}$},
בואו נבין מה תת-קבוצות של העולם צריכות לקיים\textquotedblright{} אבל
עכשיו אנחנו מתחילים את הדיון עם אותה תת-קבוצה, והיא מבחינתנו כל העולם
ואין עולם חוץ ממנה. מה שנכון הוא שהתכונה ה\textquotedblright מעניינת\textquotedblright{}
שאנחנו מחפשים עבור \L{$X$} תהיה כזו שהוא יכול \textquotedblleft לרשת\textquotedblright{}
מהמרחב המקורי \textbf{אם הוא סגור בו} ואחרת אופס. עוד מעט גם נראה
את זה קורה.

ובכן, אין מה לדבר על התכונה של להיות סגור, וגם אין מה לדבר על אקסיומת
השלמות כי אין על \L{$X$} סדר. את זה ש-\L{$f$} רציפה אנחנו כמובן
עדיין מניחים, במובן הטופולוגי של רציפות שראינו שמכליל את המקורי. מה
נשאר?
\begin{itemize}
\item \textbf{בולצאנו ויירשטראס}
\item \textbf{חסום}
\end{itemize}
שני אלו באמת יתנו לנו את מה שאנחנו צריכים, בניסוח מספיק זהיר שלהם.
אבל למעשה, אפשר להסתפק בתכונה חזקה פחות מבולצאנו ויירשטראס עצמו. בולצאנו
ויירשטראס אומר שלכל סדרה יש \textbf{תת-סדרה מתכנסת}. איך מוכיחים אותו?
ההוכחה שלו ל-\L{$\mathbb{R}$} היא ההדגמה הקלאסית לבדיחת \textquotedblleft איך
תופסים אריה במדבר?\textquotedblright{} שדיברתי עליה \href{https://gadial.net/2009/06/07/lion_in_the_desert/}{כאן}.
לוקחים את התחום שבו אברי הסדרה חיים, חוצים אותו באמצע, יש אינסוף איברים
באחד החצאים אז הולכים לחצי הזה, חוצים אותו באמצע וכן הלאה וכן הלאה.
בסוף מקבלים סדרה של קטעים שמוכלים זה בתוך זה והולכים וקטנים אל אפס;
כשבוחרים איבר אחד מכל קטע שכזה, מקבלים \textbf{סדרת קושי} ובזכות התכונות
היפות של \L{$\mathbb{R}$}, כל סדרת קושי מתכנסת והסיפור נגמר.

למה זה לא עובד ב-\L{$\mathbb{Q}$}? ראינו קודם ש-\L{$\sqrt{2}$} יכול
להיות מספר בעייתי כי הוא לא ב-\L{$\mathbb{Q}$}, בהקשר של הפונקציה
שהיא רציפה אבל עדיין אין לה מקסימום; אותו הדבר גם עם בולצאנו ויירשטראס,
אנחנו עלולים לבנות סדרת קושי שב-\L{$\mathbb{R}$} מתכנסת אל \L{$\sqrt{2}$}
אבל ב-\L{$\mathbb{Q}$} לא מתכנסת לשום דבר כי המספר היחיד שהיא יכולה
להתכנס אליו לא שייך למרחב. אז כדי שיהיה לנו סיכוי להכליל את בולצאנו
ויירשטראס אל \L{$\left(X,d\right)$} אנחנו צריכים שכל סדרת קושי ב-\L{$X$}
תתכנס. למרחב מטרי \L{$X$} שמקיים את התכונה הזו קוראים \textbf{מרחב
מטרי שלם}. השלמות הזו היא אחת משתי התכונות שאנחנו הולכים לדרוש מ-\L{$X$}.

שלמות היא תכונה של מרחבים \textbf{מטריים}, לא טופולוגיים כלליים, כי
כדי להגדיר סדרת קושי צריך מטריקה. הסדרה \L{$\left\{ a_{n}\right\} ^{\infty}_{n=1}$}
היא סדרת קושי אם לכל \L{$\varepsilon>0$} קיים \L{$N$} כך שלכל \L{$n,m>N$}
מתקיים \L{$d\left(a_{n},a_{m}\right)<\varepsilon$}. \href{https://gadial.net/2024/10/12/real_numbers_two_completness/}{בפוסט שלי}
שכבר הזכרתי פה על מושג השלמות ב-\L{$\mathbb{R}$}, התכונה הזו מכונה
\textbf{שלמות-קנטור} )על שם גאורג קנטור שביסס עליה את הבניה שלו למספרים
הממשיים(. יש הרבה דברים מעניינים לומר על התכונה הזו באופן כללי, אבל
נשמור את זה לפוסט אחר. כרגע אני רוצה אותה בשביל בולצאנו ויירשטראס.
אז בואו נעבור לדבר על התכונה השניה שאני צריך - שהמרחב יהיה \textbf{חסום}.

חסימות היא תכונה שקל לנסח במרחב מטרי: המרחב \L{$\left(X,d\right)$}
הוא חסום אם קיים \L{$a\in X$} ו-\L{$r\ge0$} כך ש-\L{$X\subseteq B\left(a,r\right)$}.
כלומר, כל המרחב כולו נכנס בתוך כדור פתוח אחד. בפרט זה אומר שהמרחק
המקסימלי בין זוג איברים במרחב הוא, אה, חסום. בבירור \L{$\mathbb{R}$}
לא מקיים את זה )אין דרך לכסות אותו עם כדור פתוח בודד, רק עם אינסוף(
ובבירור כל קטע סופי, כלומר ששני הקצוות שלו הם מספרים ממשיים, כן מקיים
את זה.

העניין הוא שזה לא מספיק. אני לא אוכל להוכיח את בולצאנו-ויירשטראס ככה.
בואו נסתכל על דוגמא נגדית טיפשית ממש: מרחב \textbf{דיסקרטי}, כלומר
קבוצה \L{$X$} עם מטריקה 

\L{$d\left(a,b\right)=\begin{cases}
1 & a\ne b\\
0 & a=b
\end{cases}$}

זו מטריקה מאוד מטופשת - המרחק בין כל שני איברים הוא {\beginL 1\endL}
חוץ מאשר המרחק שהוא {\beginL 0\endL} בין איבר לעצמו. העניין הוא שדווקא
בגלל שזו כזו סיטואציה מטופשת, היא טובה לדוגמאות נגדיות - קל להבין
אותה, והיא יכולה לשבור לא מעט דברים שמתבססים על האינטואיציה שלנו מ-\L{$\mathbb{R}$}.

במקרה הנוכחי, בוודאי שהמרחב חסום - קחו כדור ברדיוס {\beginL 1\endL}
סביב כל איבר של \L{$X$}. בנוסף, המרחב שלם בצורה מטופשת - \textbf{אין
סדרות קושי} במרחב הזה חוץ מאשר סדרות קבועות, כי בסדרת קושי עבור \L{$\varepsilon=\frac{1}{2}$}
מתישהו כל אברי הסדרה יצטרכו להיות במרחק פחות מ-\L{$\frac{1}{2}$},
כלומר כולם צריכים להיות שווים.

האם זה אומר שלכל סדרה יש תת-סדרה מתכנסת? בוודאי שלא. נגדיר \L{$X=\mathbb{N}$}
וניקח את הסדרה \L{$a_{n}=n$}, כלומר \L{$0,1,2,\ldots$}. מצד אחד,
כל סדרה מתכנסת חייבת להיות קבועה; מצד שני בסדרה שלנו אין איבר שמופיע
פעמיים, כלומר אין לה תת-סדרה קבועה. זה מסיים את הסיפור.

אוקיי... מה חסר לנו? אם נלך אל ההוכחה המקורית של בולצאנו-ויירשטראס,
ואם אנחנו חושבים על הקטע הסגור שבו ההוכחה מתעסקת בתור כל העולם, אז
השלב הראשון, שבו מחלקים את הקטע לשניים, זו בעצם הטענה \textquotedblleft לא
רק שכל העולם נכנס לתוך כדור אחד ברדיוס \L{$R$}, אלא שגם כל העולם
נכנס לתוך\textbf{ שני} כדורים ברדיוס \L{$\frac{R}{2}$}\textquotedblright .
במילים אחרות, אנחנו רוצים להיות מסוגלים להתמקד על איזורים \textbf{קטנים
יותר} של העולם, ועדיין להיות מסוגלים לכסות כל פינה אפשרית בעולם שאולי
נזדקק לה.

את זה אין במרחב דיסקרטי. כדור ברדיוס \L{$\frac{1}{2}$} מכיל בדיוק
איבר אחד - את המרכז שלו. אז אי אפשר לכסות את כל המרחב עם שניים כאלו,
ואפילו לא עם מספר \textbf{סופי} של כאלו - צריך אינסוף. וכידוע, המעבר
ממשהו סופי למשהו אינסופי הוא אולי חסר כאב, אבל הוא מביא שינויים רבים.

אפשר לחשוב על מה שקורה בבולצאנו ויירשטראס ככה: בכל שלב אנחנו לא מחלקים
את הקטע \textbf{הנוכחי} לשני חלקים, אלא מחלקים \textbf{את כל העולם}
למספר סופי כלשהו של חלקים. בגלל שזה מספר סופי, באחד מהחלקים הללו יש
אינסוף איברים מהסדרה שעדיין רלוונטיים לנו, ובחלק הזה אנחנו מתמקדים.
אם היינו נאלצים לחלק את העולם למספר \textbf{אינסופי} של חלקים, בהחלט
ייתכן שבכל אחד מהחלקים הללו היה רק מספר \textbf{סופי} של איברים, וההוכחה
הייתה הולכת לאבדון. לכן הפער הזה בין סופי ואינסופי הוא לב העניין פה
- זה, והיכולת לחלק את העולם לכדורים \textbf{קטנים כרצוננו}.

זה מוביל אותנו להגדרה הבאה, של התכונה המהותית השניה שאנחנו זקוקים
לה: נאמר ש-\L{$\left(X,d\right)$} הוא \textbf{חסום לחלוטין} אם לכל
\L{$\varepsilon>0$} קיים מספר סופי של כדורים מרדיוס \L{$\varepsilon$}
שמכסה את \L{$X$}, כלומר יש \L{$a_{1},\ldots,a_{n}\in X$} כך ש-\L{$X\subseteq\bigcup^{n}_{i=1}B\left(a_{i},\varepsilon\right)$}.

במרחב שהוא חסום לחלוטין וגם שלם קל להוכיח את בולצאנו ויירשטראס בדיוק
באופן שתיארנו: בהינתן סדרה אינסופית נבנה לה תת-סדרה מתכנסת \L{$\left\{ a_{n}\right\} ^{\infty}_{n=1}$}
בשלבים. בשלב ה-\L{$n$} נחלק את המרחב למספר סופי של כדורים ברדיוס
\L{$\frac{1}{n}$}. יש מספר סופי של כדורים, ונותרו בסדרה אינסוף איברים
שרלוונטיים לנו, אז לפחות באחד מהכדורים שנסמן ב-\L{$B_{n}$} יש אינסוף
מאברי הסדרה. נסתכל על הכדור הזה וניקח את \L{$a_{n}$} להיות אחד מאברי
הסדרה שבתוכו. כעת נשכח מכל אברי הסדרה למעט האיברים שנמצאים בתוך הכדור
שבחרנו ובאים אחרי \L{$a_{n}$} בסדרה - זה עדיין מספר אינסופי. כך נמשיך
לשלב הבא.

כדי לראות שהסדרה היא קושי, ניקח \L{$\varepsilon>0$} כלשהו ונבחר \L{$N$}
טבעי כך ש-\L{$\frac{2}{N}<\varepsilon$}. ניקח שני \L{$n,m>N$}. אז
אנחנו יודעים ש-\L{$a_{n},a_{m}$} נבחרו לסדרה \textbf{אחרי} השלב ה-\L{$N$}.
בשלב ה-\L{$N$} הותרנו בסדרה רק את מי שנמצאים בכדור \L{$B_{N}$} שרדיוסו
\L{$\frac{1}{N}$}. נסמן את מרכז הכדור ב-\L{$x$}. על פי אי שוויון
המשולש:

\L{$d\left(a_{n},a_{m}\right)\le d\left(a_{n},x\right)+d\left(a_{m},x\right)\le\frac{1}{N}+\frac{1}{N}=\frac{2}{N}<\varepsilon$}

זה מסיים את ההוכחה.

עכשיו אפשר להתחיל לתת לדברים שמות: למרחב טופולוגי שבו לכל סדרה יש
תת-סדרה מתכנסת קוראים \textbf{מרחב קומפקטי סדרתית}. שימו לב - אני
לא מדבר רק על מרחבים מטריים, כבר ראינו שאפשר לדבר על התכנסות גם במרחבים
טופולוגיים כלליים. מה שהוכחנו כרגע הוא שאם מרחב מטרי הוא שלם וחסום
לחלוטין, אז הוא קומפקטי סדרתית. עבור מרחבים מטריים הקומפקטיות הסדרתית
מספיקה כדי להוכיח את משפטי ויירשטראס; העניין הוא שבמרחב טופולוגי כללי
זה לא המצב. מבחינה היסטורית, המילה \textbf{קומפקטיות} בהקשר הטופולוגי
דיברה בדיוק על העניין הזה, של קיום תת-סדרה מתכנסת, אבל ככל שהתחום
התקדם ונהיה ברור שיש תכונה יותר חזקה שיותר מועיל להתמקד בה, המילה
\textbf{קומפקטיות} עברה אליה והתכונה הנוכחית נשארה עם שם פרס הניחומים
\textquotedblleft קומפקטיות סדרתית\textquotedblright .

אם \L{$\left(X,d\right)$} הוא קומפקטי סדרתית ו-\L{$f:X\to\mathbb{R}$}
היא רציפה, אז משפטי ויירשטראס מגיעים כמעט מעצמם, ועם אותה הוכחה שראינו
קודם: ראשית נניח ש-\L{$f$} לא חסומה ונבנה סדרה \L{$\left\{ x_{n}\right\} ^{\infty}_{n=1}$}
של אברי \L{$X$} כך ש-\L{$f\left(x_{n}\right)>n$}. נשתמש בקומפקטיות
הסדרתית של \L{$X$} כדי למצוא לה תת-סדרה שמתכנסת אל \L{$x$}, ומהרציפות
של \L{$f$} ינבע שמצד אחד \L{$f\left(x\right)$} סופי ומצד שני הוא
שווה לגבול של \L{$\lim_{n\to\infty}f\left(x_{n}\right)=\infty$} -
סתירה. עכשיו, נשתמש \textbf{באקסיומת השלמות} על \L{$\mathbb{R}$}
כדי לקבל חסם עליון \L{$M$} של התמונה של \L{$f$}, נבנה סדרה \L{$\left\{ x_{n}\right\} ^{\infty}_{n=1}$}
כך ש-\L{$f\left(x_{n}\right)\to M$}, ניקח תת-סדרה מתכנסת שלה והגבול
\L{$x$} של תת-הסדרה הזו ייתן את \L{$M$}. הקומפקטיות הסדרתית סיימה
את הסיפור יפה מאוד.

איפה השתמשנו פה בכך ש-\L{$X$} הוא מרחב מטרי? לא השתמשנו. בעצם, לא
השתמשנו כאן כמעט בכלום על \L{$X$}. הדבר היחיד שהיה רלוונטי הוא הקומפקטיות
הסדרתית. במילים אחרות, ההוכחה הזו עובדת לכל מרחב טופולוגי שהוא קומפקטי
סדרתית; פשוט עבור מרחבים מטריים היה לי קל יותר להבין בהתחלה את התכונות
\textquotedblleft מרחב שלם\textquotedblright{} ו\textquotedblright מרחב
חסום לחלוטין\textquotedblright{} שמהן הקומפקטיות הסדרתית הזו התקבלה;
התכונות הללו הן בעלות משמעות רק במרחב מטרי, אבל קומפקטיות סדרתית היא
מושג רחב יותר.

\section{חזרה לטופולוגיה כללית, לפני יומיים}

היינו יכולים לעצור כאן אם כל המוטיבציה שלנו הייתה משפטי ויירשטראס.
הבעיה היא שקשה למצוא מוטיבציה שקל להציג למה יש צורך בהגדרה כללית יותר.
קומפקטיות סדרתית שקולה להגדרה הכללית יותר בכל מרחב מטרי, ואפילו בכל
מרחב \textbf{מטריזבילי}, כלומר מרחב טופולוגי שאפשר למצוא מטריקה עליו
שתיתן את הטופולוגיה שלו. זה אומר שמלכתחילה, להראות סיטואציות שבהן
משהו עובד עם קומפקטיות אבל לא עם קומפקטיות סדרתית יצריך אותנו להבין
עולמות טופולוגיים מוזרים למדי. אז בינתיים אני לא אראה דברים שנשברים;
בתקווה אני אוכל לחזור לזה בצורה טובה יותר בפוסט נפרד שידבר על \textbf{רשתות}
)\L{Nets}(. הרעיון ברשתות הוא להכליל את מושג הסדרה, כך שהיא לא תהיה
אוסף של איברים שסדורים בסדר לינארי אלא בסדר טיפה יותר מבולגן שעדיין
יש לו כיוון ברור; בעזרת רשתות אפשר לנסח את ההגדרה הכללית של קומפקטיות
בעזרת בולצאנו ויירשטראס: מרחב הוא קומפקטי אם ורק אם לכל רשת בו יש
תת-רשת מתכנסת. אז אם נבין רשתות ובאילו סיטואציות הן עובדות למרות שסדרות
רגילות לא, יהיה יותר ברור למה צריך הגדרה כללית יותר.

לעת עתה בואו פשוט ניתן אותה. אבל כדי שזה יהיה כיף בואו ניתן \textbf{שתי}
הגדרות שקולות, אחת עם איחודים ואחת עם חיתוכים.
\begin{itemize}
\item מרחב טופולוגי \L{$X$} הוא \textbf{קומפקטי} אם לכל כיסוי פתוח שלו
יש תת-כיסוי סופי. כלומר, פורמלית: לכל אוסף \L{$\left\{ U_{\alpha}\right\} _{\alpha\in\Lambda}$}
של קבוצות פתוחות ב-\L{$X$} כך ש-\L{$\bigcup_{\alpha\in\Lambda}U_{\alpha}=X$}
קיים תת-אוסף \L{$\Lambda^{\prime}\subseteq\Lambda$} \textbf{סופי}
כך ש-\L{$\bigcup_{\alpha\in\Lambda^{\prime}}U_{\alpha}=X$}
\item מרחב טופולוגי \L{$X$} הוא \textbf{קומפקטי} אם לכל אוסף של קבוצות
סגורות שמקיים את תכונת החיתוך הסופי יש חיתוך לא ריק. כלומר, פורמלית:
לכל אוסף \L{$\left\{ C_{\alpha}\right\} _{\alpha\in\Lambda}$} של
קבוצות סגורות ב-\L{$X$} כך שלכל תת-אוסף \textbf{סופי} \L{$\Lambda^{\prime}\subseteq\Lambda$}
מתקיים \L{$\bigcap_{\alpha\in\Lambda^{\prime}}C_{\alpha}\ne\emptyset$}
מתקיים ש-\L{$\bigcap_{\alpha\in\Lambda}C_{\alpha}\ne\emptyset$}
\end{itemize}
אני חושב שטוב לראות את שתי ההגדרות הללו אחת ליד השניה מראש, פשוט כי
הדמיון ביניהן זועק לשמיים - זה שוב העניין הזה שנראה ששתי ההגדרות הן
\textbf{דואליות} - הן מאוד דומות עד כדי היפוכי תפקידים מסוימים, כאילו
אנחנו אליס ועברנו לארץ שמבעד לראי. יש לנו קבוצות \textbf{פתוחות} מול
\textbf{סגורות}; \textbf{איחודים} מול \textbf{חיתוכים}; הקבוצה \textbf{הריקה}
מול \textbf{כל} \L{$X$}; ובהגדרה הראשונה מתחילים עם זה שהכל הוא כיסוי
וממנו מגיעים \textbf{לקיום} תת אוסף סופי, ולעומת זאת בהגדרה השניה
מתחלים עם משהו של \textbf{כל} תתי האוספים הסופיים ומשם מגיעים לכך
שהכל בעל חיתוך לא ריק. עם כזו דואליות, ההוכחה ששני אלו שקולים די כותבת
את עצמה: פשוט צריך את ההגדרות בצורה מספיק פורמלית כדי שאפשר יהיה לקחת
את השקול הלוגי המתאים.

בואו נתחיל עם תזכורת מלוגיקה: אם \L{$A,B$} הן שתי טענות שכל אחת יכולה
להיות אמת או שקר, אז \L{$A\Rightarrow B$} זה \textquotedblleft אם
\L{$A$} אז \L{$B$}\textquotedblright{} - פורמלית, זו טענה שהיא שקר
רק אם \L{$A$} הוא אמת ו-\L{$B$} הוא שקר, ובכל יתר המקרים היא אמת
)ולכן \textquotedblleft אם כדור הארץ חלול אז משפט פיתגורס הוא המצאה
של האילומינטי\textquotedblright{} היא טענה נכונה לחלוטין(.עכשיו אפשר
להסתכל גם על הפסוק \L{$\neg B\Rightarrow\neg A$} )\L{$\neg$} של
משהו הוא \textbf{שלילה} - שקר הופך לאמת, אמת הופכת לשקר(. מתי הוא
מקבל ערך שקר? רק אם \L{$\neg B$} הוא אמת, כלומר אם \L{$B$} הוא שקר,
\textbf{וגם} \L{$\neg A$} הוא שקר, כלומר \L{$A$} הוא אמת. במילים
אחרות, המקרה היחיד שבו זה מקבל ערך שקר הוא שוב כאשר \L{$A$} הוא אמת
ו-\L{$B$} הוא שקר - הפסוק \L{$A\Rightarrow B$} שקול לפסוק \L{$\neg B\Rightarrow\neg A$}.
אז מה שנעשה הוא לנסח את ההגדרה לקומפקטיות עם קבוצות פתוחות \L{$A\Rightarrow B$},
ואז נקבל את זו עם קבוצות סגורות בתור \L{$\neg B\Rightarrow\neg A$}.

עבור הגדרת הקבוצות הפתוחות, \L{$A$} הוא הטענה \L{$\bigcup_{\alpha\in\Lambda}U_{\alpha}=X$}
ואילו \L{$B$} היא הטענה \textquotedblleft קיים תת-אוסף \L{$\Lambda^{\prime}\subseteq\Lambda$}
\textbf{סופי} כך ש-\L{$\bigcup_{\alpha\in\Lambda^{\prime}}U_{\alpha}=X$}\textquotedblright .
אז אני אכתוב:

\L{$\left[\bigcup_{\alpha\in\Lambda}U_{\alpha}=X\right]\Rightarrow\left[\exists\Lambda^{\prime}\left(\bigcup_{\alpha\in\Lambda^{\prime}}U_{\alpha}=X\right)\right]$}

אם רוצים להיות ממש פדנטיים, לכתוב \L{$\exists\Lambda^{\prime}$} עדיין
לא מבטיח לי ש-\L{$\Lambda^{\prime}\subseteq\Lambda$} ושהוא סופי;
כתיב יותר פורמלי יהיה \L{$\exists\Lambda^{\prime}\left(\left(\Lambda^{\prime}\subseteq\Lambda\wedge\left|\Lambda^{\prime}\right|<\infty\right)\Rightarrow\bigcup_{\alpha\in\Lambda^{\prime}}U_{\alpha}=X\right)$}.

עכשיו, בואו ניקח את השלילה. השלילה של \L{$\exists\Lambda^{\prime}\left(\bigcup_{\alpha\in\Lambda^{\prime}}U_{\alpha}=X\right)$}
היא \L{$\forall\Lambda^{\prime}\left(\bigcup_{\alpha\in\Lambda^{\prime}}U_{\alpha}\ne X\right)$}.
עכשיו, מה זה אומר ש-\L{$\bigcup_{\alpha\in\Lambda^{\prime}}U_{\alpha}\ne X$}?
זה אומר שאם ניקח משלימים לשני האגפים, אז \L{$\left(\bigcup_{\alpha\in\Lambda^{\prime}}U_{\alpha}\right)^{c}\ne\emptyset$},
כלומר על פי כללי דה-מורגן \L{$\bigcap_{\alpha\in\Lambda^{\prime}}U^{c}_{\alpha}\ne\emptyset$}.
נסמן \L{$C_{\alpha}=U^{c}_{\alpha}$} ומכיוון ש-\L{$U_{\alpha}$}
פתוחה נקבל ש-\L{$C_{\alpha}$} סגורה, וזה בדיוק מה שרצינו. באופן דומה
השלילה של \L{$\bigcup_{\alpha\in\Lambda}U_{\alpha}=X$} היא \L{$\bigcap_{\alpha\in\Lambda}C_{\alpha}\ne\emptyset$},
מה שמסיים את ההוכחה.

אוקיי, אז יש לנו שתי הגדרות מוזרות, אבל איך הן קשורות למה שדיברנו
עליו קודם? בתור התחלה, איך אני מוכיח את משפטי ויירשטראס עם קומפקטיות
ולא עם קומפקטיות סדרתית? זה נכון שקומפקטיות סדרתית מספיקה כדי להוכיח
את ויירשטראס, אבל - זה שמרחב הוא קומפקטי \textbf{לא מבטיח} שהוא יהיה
קומפקטי סדרתית; כמו שאמרתי למעלה, זה עובד במרחבים מטריזביליים אבל
באופן כללי בשביל שזה יעבוד צריך רשתות ואני לא מדבר עליהן כאן.

אז יש לנו סיטואציה מעניינת - אנחנו רוצים להוכיח את משפט ויירשטראס
בשיטה שתצטרך להיות שונה באופיה ממה שאנחנו מכירים מאינפי. מה שנחמד
הוא שזו גם שיטה פשוטה מאוד. אבל קודם, בואו ננסח את המשפט למרחבים טופולוגיים
בצורה הכי כללית שאנחנו יכולים. בשביל זה, אני יכול שהפונקציה לא תחזיר
פלט שהוא \textbf{מספר ממשי}; יהיה לי מספיק שהיא תלך למרחב שהוא \textbf{קבוצה
סדורה לינארית} עם טופולוגיית הסדר.

אז פורמלית, תהא \L{$f:X\to Y$} רציפה כך ש-\L{$X$} מרחב קומפקטי ו-\L{$Y$}
מרחב עם טופולוגיית הסדר, ונוכיח שקיימים \L{$m,M\in X$} כך ש-\L{$f\left(m\right)\le f\left(x\right)\le f\left(M\right)$}
לכל \L{$x\in X$} )עם יחס הסדר שמוגדר על \L{$Y$}(.

הדבר הראשון שאזדקק לו הוא טענה מאוד שימושית בפני עצמה. אם \L{$f:X\to Y$}
היא פונקציה \textbf{רציפה} ו-\L{$X$} קומפקטי, אז \L{$f\left(X\right)$}
קומפקטית. כלומר, התמונה של קבוצה קומפקטית על ידי פונקציה רציפה גם
היא קומפקטיות - קומפקטיות \textbf{משתמרת} תחת פונקציות רציפות. הנה
הוכחה ישירות על פי ההגדרה. נסמן \L{$A=f\left(X\right)$} ונחשוב על
\L{$A$} בתור תת-מרחב של \L{$Y$} ביחס לטופולוגיית תת-המרחב. ניקח
כיסוי פתוח \L{$\bigcup_{\alpha}V_{\alpha}=A$} ונגדיר \L{$U_{\alpha}=f^{-1}\left(V_{\alpha}\right)$}.
אלו קבוצות פתוחות כי \L{$f$} רציפה; ולכל \L{$x\in X$} מתקיים ש-\L{$f\left(x\right)\in Y^{\prime}$},
כלומר קיים \L{$V_{\alpha}$} כך ש-\L{$f\left(x\right)\in V_{\alpha}$}
ולכן \L{$x\in f^{-1}\left(V_{\alpha}\right)=U_{\alpha}$}, כך שה-\L{$U_{\alpha}$}
הללו מהוות כיסוי פתוח של \L{$X$} ולכן יש תת-כיסוי סופי \L{$U_{1},\ldots,U_{n}$}.
עכשיו, \L{$V_{1},\ldots,V_{n}$} יהיה תת-כיסוי סופי של \L{$A$}, מה
שמסיים את ההוכחה. 

עכשיו אפשר לעבור להוכחת משפט ויירשטראס עצמו. בואו כמו קודם נסמן \L{$A=f\left(X\right)$},
אז אני יודע שזו קבוצה קומפקטית - כלומר, \L{$A$} הוא מרחב קומפקטי
)ביחס לטופולוגיה שהוא יורש מ-\L{$Y$}, דהיינו טופולוגיית תת-המרחב(.
מה שאני בעצם רוצה להראות הוא שקיים ב-\L{$A$} איבר מקסימלי )ומינימלי,
אבל ההוכחה זה אותו דבר( - לאיבר הזה יש מקור תחת \L{$f$} כי לכל איבר
של \L{$A$} יש מקור תחת \L{$f$}, אז ניקח את \L{$M$} להיות אחד מהמקורות.
כלומר: הפואנטה של משפט ויירשטראס )וזה נכון גם למשפט המקורי( היא שלקבוצה
קומפקטית \textbf{סדורה} יש מקסימום ומינימום.

כדי לראות את זה בפעולה, בואו נסתכל על כל הקבוצות מהצורה \L{$\left(-\infty,a\right)\cap A$}
כך ש-\L{$a\in A$} - כלומר לקחנו את הקרניים האינסופיות לשמאל ב-\L{$Y$}
עם נקודת קצה מתוך \L{$A$}, ואז הצמצמנו למרחב \L{$A$}; לפי טופולוגיית
תת-המרחב, הקבוצות הללו הן פתוחות ב-\L{$A$}. עכשיו, אם ב-\L{$A$}
\textbf{אין} איבר גדול ביותר, הקבוצות הללו מכסות את כל \L{$A$}, כי
לכל \L{$a_{1}\in A$} קיים \L{$a_{2}\in A$} כך ש-\L{$a_{1}<a_{2}$}
ולכן \L{$a_{1}\in\left(-\infty,a_{2}\right)\cap A$}. אז אם \textbf{נניח
בשלילה} שאין ב-\L{$A$} איבר מקסימלי, נקבל כיסוי פתוח של \L{$A$}
ובגלל הקומפקטיות של \L{$A$} אפשר לקחת לו תת-כיסוי סופי שהוא מהצורה
\L{$\left(-\infty,a_{1}\right)\cap A,\ldots,\left(-\infty,a_{n}\right)\cap A$}.
בואו נניח שנקודות הקצה מסודרות על פי יחס הסדר, כלומר \L{$a_{1}<a_{2}<\ldots<a_{n}$},
אז מצד אחד, \L{$a_{n}\in A$} כי לקחנו מלכתחילה רק קבוצות שבהן נקודת
הקצה שייכת ל-\L{$A$}; ומצד שני האיחוד של כל הקבוצות הללו מכיל את
כל \L{$A$} אבל הוא \textbf{לא} מכיל את \L{$a_{n}$} כי \L{$a_{n}$}
גדול מכל הנקודות בקטעים הללו. זו סתירה שמראה שההנחה שב-\L{$A$} אין
איבר גדול ביותר הייתה שגויה, וזה מסיים את ההוכחה. וואו זה היה פשוט.
שום בנייה של סדרות ובולצאנו-ויירשטראס; רק טיעון מאוד פשוט על כיסויים.
זה די מעניין כמה הגדרה כל כך מוזרה היא חזקה בפועל.

\section{ההווה, כלומר - לאן עכשיו?}

כן, זה אני, אתם בוודאי מבינים איך הגעתי למצב הזה, אבל עכשיו השאלה
היא לאן מתקדמים מכאן. יש די הרבה דברים להגיד על קומפקטיות אבל אני
לא אעשה רשימת מכולת. יש יעד קונקרטי שאני רוצה להציג: \textbf{משפט
טיכונוף}, המשפט שמראה שמכפלה של מרחבים קומפקטיים היא קומפקטית. זה
אחד מהמשפטים ה\textquotedblright גדולים\textquotedblright{} בקורס מבוא
לנושא ואני אקדיש פוסט שלם להוכחה שלו; אבל אני רוצה לתת עכשיו קצת מוטיבציה
לכמה חזק המשפט הזה יכול להיות ובאיזו צורה מוזרה, על ידי כך שאני אכניס
לתמונה תוצאה מתחום שנראה לא קשור בעלילה - משפט \textbf{מלוגיקה מתמטית}
שנקרא, כמה הולם, \textquotedblleft משפט הקומפקטיות\textquotedblright .
תיארתי את הסיפור המרהיב הזה \href{https://gadial.net/2008/03/03/compactness_theorem_intro/}{בראשית ימי הבלוג},
אבל בואו נעשה את זה שוב בשולי הפוסט הזה.

מה שאני מדבר עליו פה הוא החלק של לוגיקה שנקרא \textbf{תחשיב הפסוקים}.
בדרך כלל בלוגיקה אוהבים לדבר על \textbf{לוגיקה מסדר ראשון} וגם שם
יש משפט קומפקטיות, אבל הוא מסובך יותר והחלק הטופולוגי לא מכסה בו את
הכל \href{https://gadial.net/2013/02/25/godel_completeness_proof_1/}{וכבר יש לי פוסט}
על מה שנקרא משפט השלמות של גדל שממנו אפשר להסיק את משפט הקומפקטיות
ללוגיקה מסדר ראשון, אז כאן נתמקד בתחשיב הפסוקים שהוא מין גרסת חימום
שכזו שכוללת מספיק רעיונות יפים בפני עצמה.

כבר ראינו את תחשיב הפסוקים בפוסט הזה, כשדיברתי על טענה מהצורה \L{$A\Rightarrow B$}.
באופן כללי כל דבר בתחשיב הפסוקים בנוי ככה: יש לנו \textbf{משתנים}
שיכולים לקבל ערך \textquotedblleft שקר\textquotedblright{} או \textquotedblleft אמת\textquotedblright{}
שאני פשוט אסמן ב-{\beginL 0\endL} ו-{\beginL 1\endL} מעכשיו והלאה;
ויש לנו \textbf{פסוקים} שנבנים מתוך המשתנים על ידי חיבור שלהם עם \textbf{קשרים
לוגיים} שהם אופרטורים על ערכי אמת. את \L{$\Rightarrow$} כבר ראינו,
אם כי כשכותבים פסוקים פורמלית מעדיפים להשתמש ב-\L{$\rightarrow$}
ולא בחץ כפול. ויש גם את \L{$\wedge$} עבור \textquotedblleft וגם\textquotedblright{}
ואת \L{$\vee$} עבור \textquotedblleft או\textquotedblright{} ואת \L{$\neg$}
עבור שלילה - ואנחנו לא באמת צריכים להבין עד הסוף אף אחד מהם כדי להבין
את משפט הקומפקטיות. רק שמשתמשים בהם לבנות פסוקים ושבנוסף לכך, \textbf{וזה
לב העניין פה}, הבניה הזו היא תמיד \textbf{סופית} - אחרי מספר סופי
של צעדים אנחנו עוצרים ומקבלים את הפסוק שלנו, שהוא סדרה סופית של סימבולית,
ובפרט יש בו גם מספר סופי של משתנים.

אנחנו אומרים שפסוק הוא \textbf{ספיק} אם יש הצבה של ערכים למשתנים שתיתן
לפסוק כולו את הערך {\beginL 1\endL} )\textquotedblleft אמת\textquotedblright (.
בהרבה סיטואציות, האתגר האמיתי הוא למצוא השמה שמספקת בו זמנית הרבה
פסוקים שונים, אם קיימת כזו. למשל, דמיינו לוח סודוקו. זה לוח \L{$9\times9$}
שבו כל משבצת יכולה להכיל מספר בין {\beginL 1\endL} ל-{\beginL 9\endL}.
אני יכול \textbf{למדל} את המילוי של הלוח באמצעות משתנים: יהיה לי משתנה
מהצורה \L{$X^{a}_{i,j}$}שאומר \textquotedblleft בתא \L{$\left(i,j\right)$}
יש את הערך \L{$a$}\textquotedblright{} אם הוא מקבל {\beginL 1\endL}
וההפך אם הוא מקבל {\beginL 0\endL}. בסך הכל יש \L{$9\times9\times9$}
משתנים כאלו כי \L{$1\le i,j\le9$} כי זה הגודל של הלוח ו-\L{$1\le a\le9$}
כי אלו הערכים שאנחנו שמים במשבצות. אלו רק המשתנים; עכשיו אני משתמש
בפסוקים כדי \textquotedblleft להכריח\textquotedblright{} את ההשמה המספקת
להתאים למילוי חוקי של לוח סודוקו. למשל, כדי להבטיח שבתא \L{$3,5$}
לא יופיעו גם {\beginL 3\endL} וגם {\beginL 7\endL} בו זמנית אני יכול
להוסיף את הפסוק \L{$\neg\left(X^{3}_{3,5}\wedge X^{7}_{3,5}\right)$},
אז פשוט לכל \L{$1\le i,j\le9$} וכל \L{$1\le a\ne b\le9$} אני אוסיף
את הפסוק \L{$\neg\left(X^{a}_{i,j}\wedge X^{b}_{i,j}\right)$}. אני
גם רוצה לדרוש שבכל תא יהיה ערך אחת, כלומר לא כל המשתנים שמתאמים לו
מקבלים {\beginL 0\endL}, ואת זה עושה הפסוק \L{$\left(X^{1}_{i,j}\vee\ldots\vee X^{9}_{i,j}\right)$}.
לבסוף, בלוח סודוקו חוקי יש הרבה זוגות של משבצות שונות שאסור להן לקבל
את אותו ערך )למשל משבצות באותה שורה(. אם \L{$\left(i_{1},j_{1}\right),\left(i_{2},j_{2}\right)$}
הן משבצות כאלו, אני אוסיף את הפסוק \L{$\neg\left(X^{a}_{i_{1},j_{1}}\wedge X^{a}_{i_{2},j_{2}}\right)$}
לכל \L{$1\le a\le9$}, וזהו; מעכשיו כל מה שמספק את כל הפסוקים הללו
בו זמנית מקודד מילוי חוקי של לוח סודוקו )ואני יכול למשל להוסיף משתנה
כמו \L{$X^{7}_{5,4}$} בתור פסוק נפרד כדי לדרוש שבתא \L{$\left(5,4\right)$}
יהיה {\beginL 7\endL}, כלומר כדי לדרוש שפתרון הסודוקו שאני מוצא יתאים
לפתרון חלקי שכבר יש לי, כמו שחידות סודוקו עובדות(.

העניין הוא שאין ממש הבדל בין פסוק יחיד ובין קבוצה של פסוקים אם זו
קבוצה \textbf{סופית}, כי כל קבוצת פסוקים סופית אפשר להפוך לפסוק יחיד
על ידי \L{$\wedge$} לכל הפסוקים בקבוצה. אבל בקבוצה \textbf{אינסופית}
של פסוקים אי אפשר לנקוט בתעלול הזה כי פסוק לא יכול להיות אינסופי באורכו.
ומה שמשפט הקומפקטיות בא לומר הוא שכדי לדעת אם קבוצת פסוקים אינסופית
היא ספיקה, מספיק להצטמצם ולבדוק רק אם כל תת-קבוצת פסוקים \textbf{סופית}
שלה היא ספיקה. זו תוצאה מאוד מפתיעה במבט ראשון, כי גם אם אפשר לספק
כל תת-קבוצת פסוקים סופית, איך בדיוק בונים מזה השמה מספקת אחת שעובדת
עבור כל אינסוף הפסוקים בו זמנית? בפרט שזה יכול להיות אפילו אינסוף
לא בן מניה? ההשמות המספקות עבור הקבוצות הסופיות עשויות מאוד לא להסכים
זו עם זו, איך אפשר לשלב אותן להשמה אחת? ובכן, קומפקטיות. וכשנוכיח
את משפט טיכונוף נראה ששם צצות אותן שאלות קשות שהעליתי כאן, רק שבהוכחה
של טיכונוף אין דרך לטאטא אותן מתחת לשטיח עם להטוטים כמו אלו שתכף אעשה;
משפט טיכונוף זה המקום שבו תופסים את השור בקרניו ומתמודדים איתו בצורה
מאוד יפה.

הנה הניסוח הפורמלי:
\begin{itemize}
\item )משפט הקומפקטיות לתחשיב הפסוקים(: תהא \L{$\Phi$} קבוצת פסוקים. אם
כל תת-קבוצה \textbf{סופית} \L{$\Phi^{\prime}\subseteq\Phi$} היא ספיקה,
אז \L{$\Phi$} ספיקה.
\end{itemize}
מזכיר משהו? כמובן, זה מאוד מזכיר קומפקטיות בניסוח הקבוצות הסגורות
שלה:
\begin{itemize}
\item יהא \L{$\Psi$} אוסף של קבוצות סגורות במרחב קומפקטי \L{$X$}. אם כל
תת-קבוצה \textbf{סופית} \L{$\Psi^{\prime}\subseteq\Psi$} מקיימת \L{$\bigcap\Psi^{\prime}\ne\emptyset$}
אז \L{$\bigcap\Psi\ne\emptyset$}
\end{itemize}
כשרואים את זה ככה, ההוכחה די כותבת את עצמה. \textquotedblleft\L{$\Phi$}
ספיקה\textquotedblright{} אומר ש\textbf{קיימת} השמה שמספקת את כל אברי
\L{$\Phi$} בו זמנית, כלומר שאם ניקח לכל פסוק ב-\L{$\Phi$} את קבוצת
ההשמות המספקות שלו \textbf{ונחתוך} את כולן, נקבל חיתוך לא ריק. אז
בואו נעשה את זה: נתאים לכל \L{$\varphi\in\Phi$} קבוצה \L{$C_{\varphi}$}
של ההשמות שמספקות את \L{$\varphi$} ועכשיו נקבל \L{$\Psi=\left\{ C_{\varphi}\ |\ \varphi\in\Phi\right\} $}
שמקיימת בדיוק ש-\L{$\bigcap\Psi\ne\emptyset$} אם ורק אם יש השמה שמספקת
בו זמנית את כל אברי \L{$\Phi$}. מה נשאר? להבין באיזה מרחב אנחנו חיים,
בעצם, למה הוא קומפקטי, ולמה כל \L{$C_{\varphi}$} שכזו היא \textbf{סגורה}.

נניח שקבוצת המשתנים שלנו היא \L{$V$}, אז \textbf{השמה} היא פונקציה
\L{$\tau:V\to\left\{ 0,1\right\} $}. כלומר, אם אני רוצה לחשוב על
מרחב טופולוגי שהשמות הן האיברים שלו, הדרך הטבעית לחשוב עליו היא בתור
מרחב \textbf{מכפלה} שמתקבל מזה שלוקחים את המרחב \L{$\left\{ 0,1\right\} $}
וכופלים אותו בעצמו \L{$\left|V\right|$} פעמים )כלומר, פורמלית, מסתכלים
על מרחב המכפלה \L{$\prod_{v\in V}\left\{ 0,1\right\} $}(. אפשר למשל
לדמיין ש-\L{$V=\left\{ X_{1},X_{2},X_{3},\ldots\right\} $} ואז מרחב
ההשמות יהיה \L{$\left\{ 0,1\right\} \times\left\{ 0,1\right\} \times\left\{ 0,1\right\} \times\ldots$},
כלומר סדרות אינסופיות של {\beginL 0\endL} ו-{\beginL 1\endL}; אבל
\L{$V$} יכולה להיות גם קבוצה גדולה יותר, לא בת מניה, וגם במקרה הזה
משפט הקומפקטיות עדיין עובד.

על המרחב \L{$\left\{ 0,1\right\} $} נשים את הטופולוגיה הדיסקרטית,
שזה אומר בפועל \textbf{שכל} תתי הקבוצות של המרחב הן פתוחות. אז במרחב
כל ההשמות \L{$\left\{ 0,1\right\} ^{V}$} עם טופולוגיית המכפלה, מה
הקבוצות הפתוחות? כזכור, תת-הבסיס לטופולוגיה הזו עובד ככה: בוחרים איבר
אחד \L{$x\in V$} ותת קבוצה \L{$B\subseteq\left\{ 0,1\right\} $}.
עכשיו מסתכלים על קבוצת כל הפונקציות \L{$\tau:V\to\left\{ 0,1\right\} $}
כך שאין לנו שום מגבלה על \L{$\tau$} בשום מקום \textbf{חוץ מאשר} עבור
הערך \L{$x$}, ושם הדרישה שלנו היא \L{$\tau\left(x\right)\in B$}.
בגלל ש-\L{$\left\{ 0,1\right\} $} היא כזו קבוצה קטנה, המגבלה היחידה
בעלת המשמעות היא לתת ל-\L{$\tau$} ערך קונקרטי. אז פורמלית, כל אברי
תת-הבסיס לטופולוגיה של \L{$\left\{ 0,1\right\} ^{V}$} הם \textquotedblleft הפונקציות
שעל \L{$x$} נותנות {\beginL 0\endL}\textquotedblright{} ו\textquotedblright הפונקציות
שעל \L{$x$} נותנות {\beginL 1\endL}\textquotedblright{} כש-\L{$x$}
הוא איבר כלשהו של \L{$V$}.

עכשיו, מה הקבוצה המשלימה לקבוצת כל הפונקציות שעל \L{$x$} מחזירות
{\beginL 0\endL}? כמובן, כל הפונקציות שעל \L{$x$} מחזירות לא-{\beginL 0\endL},
כלומר מחזירות {\beginL 1\endL}. ואלו... עדיין אברי תת-הבסיס! עכשיו,
קבוצה פתוחה כללית מתקבלת מחיתוכים סופיים ואיחודים כלליים של אברי תת-הבסיס,
כך שקבוצה \textbf{סגורה} כללית תתקבל ממשלים של זה, ומשלים הופך את
התפקידים של חיתוך ואיחוד, כך שקבוצה סגורה כללית מתקבלת מאיחודים סופיים
וחיתוכים כלשהם של אברי הבסיס. מה זה חיתוך של אברי בסיס? זה להוסיף
\textbf{עוד} משתנים שעליהם הערך של השמה נקבע באופן יחיד. ומה זה איחוד
של אברי בסיס? זה להוסיף עוד השמות שהאופן שבו הן נקבעות באופן יחיד
יכול להיות שונה לגמרי.

והנה הפאנץ'. ניקח פסוק \L{$\varphi$} ונסתכל על \L{$C_{\varphi}$},
קבוצת כל ההשמות שמספקת אותו. ב-\L{$\varphi$} יש רק \textbf{מספר סופי}
של משתנים, ולכן יש רק מספר סופי של ערכים שונים שההשמות המספקות של
\L{$\varphi$} יכולות לקבל. לכל השמה קונקרטית למשתנים של \L{$\varphi$}
שמספקת אותו, הקבוצה של כל ההשמות שעל המשתנים של \L{$\varphi$} מקבלות
את הערכים הקונקרטיים הללו היא קבוצה סגורה; ואיחוד סופי של קבוצות כאלו
הוא קבוצה סגורה, והוא נותן את כל \L{$C_{\varphi}$}, כך ש-\L{$C_{\varphi}$}
היא קבוצה סגורה.

אז אם נסתכל על ה\textquotedblright מה נשאר?\textquotedblright{} שכתבתי
קודם - אנחנו מבינים כבר באיזה מרחב אנחנו חיים ולמה \L{$C_{\varphi}$}
היא קבוצה סגורה בו, אבל למה הוא קומפקטי? המרחב הדיסקרטי \L{$\left\{ 0,1\right\} $}
הוא קומפקטי בצורה טריוויאלית; כל מרחב סופי הוא קומפקטי כי מן הסתם
אם נתון לנו כיסוי, הכיסוי הזה הוא מראש סופי ואין אתגר בלמצוא לו תת-כיסוי
סופי. אז הדבר היחיד שנשאר הוא להבין למה אם לוקחים אוסף של מרחבים קומפקטיים
ומסתכלים על מרחב המכפלה שלהם בטופולוגיית המכפלה, גם המרחב הזה הוא
קומפקטי. זה \textbf{משפט טיכונוף} ולהוכחה שלו יוקדש הפוסט הבא.
\end{document}
